\chapter{Production-Grade RAG: Ingestion, Chunking, and Contextual Retrieval}
\label{chap:rag}

\begin{chaptergoalsbox}
\textbf{What You Will Learn:}
\begin{itemize}[leftmargin=15pt]
  \item Why naive character splitting corrupts tables, multi-column PDFs, and pronouns---and where each failure occurs in the ingestion pipeline.
  \item High-fidelity document parsing: coordinate-stream pitfalls, vision layout analysis, and row-to-proposition table flattening.
  \item Chunking mechanics: fixed-size, recursive separator descent (with true recursion and separator-aware overlap), and semantic cosine splitting.
  \item Transformer tokenization, attention dilution, Rotary Position Embeddings (RoPE), and the ``Lost in the Middle'' U-shaped curve.
  \item Small-to-big retrieval (\texttt{ParentDocumentRetriever}) and Anthropic's Contextual Retrieval for eliminating orphan chunks.
\end{itemize}

\textbf{What You Should Already Understand:}
\begin{itemize}[leftmargin=15pt]
  \item Vectors, dot products, and cosine similarity (Chapter~1).
  \item Python strings, regular expressions, lists, and dictionaries.
\end{itemize}

\textbf{Why This Topic Exists:}
Most Retrieval-Augmented Generation (RAG) prototypes split a document every 500 characters, embed the slices, and paste the top 5 matches into an LLM prompt. In production, that simplicity collapses: PDF columns interleave across pages, tables get chopped away from their column headers, clauses lose their subject names, and relevant facts placed in the middle of a long prompt are quietly ignored by the model.

\textbf{Motivating Practical Problem:}
A legal intelligence platform indexes 200-page corporate bond indentures. When an analyst asks, \emph{``What is the debt incurrence threshold under the negative covenants if consolidated EBITDA falls below \$50M?''}, naive RAG retrieves a fragment reading: \emph{``...provided, however, that the ratio shall not exceed 3.5:1.0 in such event.''} Because the covenant title and parent definition were sliced into a different chunk, the LLM cannot tell which clause applies and invents an incorrect legal answer.
\end{chaptergoalsbox}

\section{Anatomy of a RAG Failure}

Before writing a single line of ingestion code, trace how an answer goes wrong. Every stage between a raw enterprise file and the LLM's final response can silently corrupt meaning. Figure~\ref{fig:v2-ch02-rag-failure-map} maps the production pipeline alongside the concrete failure mode that strikes each stage when naive defaults are used.

\begin{figure}[htbp]
\centering
\begin{tikzpicture}[
  stage/.style={draw=headingblue, fill=codebg, rounded corners=2pt, font=\sffamily\scriptsize\bfseries, minimum width=2.4cm, minimum height=1.0cm, align=center},
  fail/.style={draw=warningRed, fill=warningBg, rounded corners=2pt, font=\sffamily\tiny, minimum width=2.4cm, minimum height=1.1cm, align=center},
  fix/.style={draw=intuitionTeal, fill=intuitionBg, rounded corners=2pt, font=\sffamily\tiny, minimum width=2.4cm, minimum height=1.0cm, align=center},
  arrow/.style={-{Stealth[scale=0.8]}, thick, color=linkblue},
  farrow/.style={-{Stealth[scale=0.7]}, dashed, color=warningRed}
]

\node[stage] (s1) at (0,0) {1. Layout\\Parsing};
\node[stage] (s2) at (3.0,0) {2. Boundary\\Chunking};
\node[stage] (s3) at (6.0,0) {3. Contextual\\Enrichment};
\node[stage] (s4) at (9.0,0) {4. Dual-Store\\Indexing};
\node[stage] (s5) at (12.0,0) {5. Context\\Assembly};

\draw[arrow] (s1) -- (s2);
\draw[arrow] (s2) -- (s3);
\draw[arrow] (s3) -- (s4);
\draw[arrow] (s4) -- (s5);

\node[fail] (f1) at (0,-1.55) {\textbf{Failure:} Interleaved\\2-col text;\\shredded tables};
\node[fail] (f2) at (3.0,-1.55) {\textbf{Failure:} Sentences\\cut mid-clause;\\lost overlap};
\node[fail] (f3) at (6.0,-1.55) {\textbf{Failure:} Pronoun\\``the company''\\unsearchable};
\node[fail] (f4) at (9.0,-1.55) {\textbf{Failure:} Small chunks\\lack context; large\\chunks blur vectors};
\node[fail] (f5) at (12.0,-1.55) {\textbf{Failure:} Middle\\chunks ignored\\by LLM attention};

\node[fix] (x1) at (0,-3.05) {\textbf{Fix:} Vision layout\\boxes + row\\propositions};
\node[fix] (x2) at (3.0,-3.05) {\textbf{Fix:} Recursive +\\semantic cosine\\boundaries};
\node[fix] (x3) at (6.0,-3.05) {\textbf{Fix:} Prepend doc\\+ section summary\\before embedding};
\node[fix] (x4) at (9.0,-3.05) {\textbf{Fix:} Child vectors\\point to parent\\doc store};
\node[fix] (x5) at (12.0,-3.05) {\textbf{Fix:} Rerank + place\\top hits at edges\\of prompt};

\draw[farrow] (s1) -- (f1); \draw[arrow] (f1) -- (x1);
\draw[farrow] (s2) -- (f2); \draw[arrow] (f2) -- (x2);
\draw[farrow] (s3) -- (f3); \draw[arrow] (f3) -- (x3);
\draw[farrow] (s4) -- (f4); \draw[arrow] (f4) -- (x4);
\draw[farrow] (s5) -- (f5); \draw[arrow] (f5) -- (x5);

\end{tikzpicture}
\caption{Production RAG Data Flow, Stage-Specific Failure Modes, and Architectural Mitigations.}
\label{fig:v2-ch02-rag-failure-map}
\end{figure}

Notice the core lesson of Figure~\ref{fig:v2-ch02-rag-failure-map}: if Stages 1--3 shred a table or strip away an entity name, Stage 5 cannot guess the missing facts later. High-accuracy RAG is won at ingestion time.

\section{Document Parsing and Structural Extraction}

\subsection{Why Raw PDF Text Streams Fail}
PDF is a print-drawing format, not a structured text format. Internally, a PDF stores painting commands placed at $(x, y)$ page coordinates:
\begin{equation}
(x=72,\; y=640,\; \text{``Consolidated''},\; \text{Helvetica},\; 10\text{pt})
\end{equation}
A PDF has no built-in tags for ``paragraph,'' ``column,'' or ``table.'' When a naive text extractor walks the stream left-to-right across a two-column financial report, it stitches Column A's line directly into Column B's line:
\begin{center}
\emph{``In 2024 our European division grew 14\%''} $\longleftrightarrow$ \emph{``While regulatory filings remained pending.''}
\end{center}

\begin{underthehoodbox}[Vision-Based Layout Analysis]
Modern document parsers (such as Marker, Docling, or LayoutLMv3) treat pages visually before reading strings:
\begin{enumerate}
  \item Render each PDF page into a bitmap image (e.g., 150--300 DPI).
  \item Run an object-detection model to draw bounding boxes labeled \texttt{Heading}, \texttt{Paragraph}, \texttt{Table}, \texttt{Header}, and \texttt{Footer}.
  \item Sort bounding boxes into true reading order (top-to-bottom within Column 1, then Column 2), drop repeated page headers/footers, and extract text inside each box via native PDF glyphs or OCR.
  \item Emit a clean Markdown hierarchy preserving headings and tables.
\end{enumerate}
\end{underthehoodbox}

\subsection{Flattening Tables into Self-Contained Propositions}
If a raw pipe-delimited Markdown table spans 40 rows and gets split at row 15, rows 16--40 lose their header row (\texttt{Year | Revenue | OpEx | EBITDA}) and become meaningless numbers. A reliable solution converts each row into a self-contained sentence (Listing~\ref{lst:v2-ch02-tables}).

\begin{lstlisting}[language=Python, label={lst:v2-ch02-tables}, caption={Flattening Financial Table Rows into Self-Contained Propositions}]
def table_to_context_tuples(table_data: list[dict[str, str]]) -> list[str]:
    """Flattens structured table rows into standalone sentences."""
    propositions: list[str] = []
    for row in table_data:
        prop = (
            f"In financial year {row['Year']}, the company recorded revenue "
            f"of {row['Revenue']}, operating expenses of {row['OpEx']}, "
            f"yielding an EBITDA of {row['EBITDA']}."
        )
        propositions.append(prop)
    return propositions

# Trace on a 2-row sample:
rows = [
    {"Year": "2023", "Revenue": "$120M", "OpEx": "$80M", "EBITDA": "$40M"},
    {"Year": "2024", "Revenue": "$150M", "OpEx": "$95M", "EBITDA": "$55M"},
]
# Output[0]: "In financial year 2023, the company recorded revenue of $120M,
#             operating expenses of $80M, yielding an EBITDA of $40M."
\end{lstlisting}

Each emitted string now carries its own column labels, so even if two rows land in separate chunks, neither loses its meaning.

\section{Chunking Topologies: Fixed, Recursive, and Semantic}

\subsection{Tokens vs Characters: A Quick Primer}
Before measuring chunk sizes, clarify what LLMs actually count. Language models do not read raw characters or whole English words; a \textbf{tokenizer} (typically Byte-Pair Encoding, or BPE) slices text into common subword pieces called \textbf{tokens}. For example, \texttt{"unbelievable"} might become \texttt{["un", "believ", "able"]} (3 tokens). As a rule of thumb:
\begin{itemize}[leftmargin=18pt]
  \item In English prose, $1\text{ token} \approx 4\text{ characters}$.
  \item In source code or JSON (heavy indentation, symbols, camelCase), $1\text{ token} \approx 2.5\text{ to }3.0\text{ characters}$.
\end{itemize}

\begin{figure}[htbp]
\centering
\begin{tikzpicture}[
  chunk/.style={draw=headingblue, fill=codebg, rounded corners=1pt, minimum height=0.6cm, font=\sffamily\tiny, align=center},
  sep/.style={draw=warningRed, dashed, thick}
]

\node[anchor=west, font=\sffamily\bfseries\footnotesize] at (0, 2.2) {A. Fixed-Size Chunking (500 tokens)};
\draw[chunk, minimum width=3.5cm] (0, 1.5) -- ++(3.5,0) node[midway] {Chunk 1 (500 tokens)};
\draw[chunk, minimum width=3.5cm] (3.8, 1.5) -- ++(3.5,0) node[midway] {Chunk 2 (500 tokens)};
\draw[chunk, minimum width=3.5cm] (7.6, 1.5) -- ++(3.5,0) node[midway] {Chunk 3 (500 tokens)};
\draw[sep] (3.65, 1.1) -- (3.65, 1.9) node[above, font=\tiny\sffamily, color=warningRed] {Split mid-sentence};

\node[anchor=west, font=\sffamily\bfseries\footnotesize] at (0, 0.7) {B. Recursive Separator Chunking};
\draw[chunk, minimum width=3.0cm, fill=intuitionBg] (0, 0) -- ++(3.0,0) node[midway] {Section 1.1 ($\backslash$n$\backslash$n)};
\draw[chunk, minimum width=4.5cm, fill=intuitionBg] (3.3, 0) -- ++(4.5,0) node[midway] {Section 1.2 Paragraph ($\backslash$n$\backslash$n)};
\draw[chunk, minimum width=3.2cm, fill=intuitionBg] (8.1, 0) -- ++(3.2,0) node[midway] {Section 1.3 ($\backslash$n$\backslash$n)};

\node[anchor=west, font=\sffamily\bfseries\footnotesize] at (0, -0.8) {C. Semantic Distance Chunking};
\draw[chunk, minimum width=2.4cm, fill=exampleBg] (0, -1.5) -- ++(2.4,0) node[midway] {Topic A (High Sim)};
\draw[chunk, minimum width=5.2cm, fill=exampleBg] (2.7, -1.5) -- ++(5.2,0) node[midway] {Topic B (Cosine Shift Detected)};
\draw[chunk, minimum width=3.0cm, fill=exampleBg] (8.2, -1.5) -- ++(3.0,0) node[midway] {Topic C};

\end{tikzpicture}
\caption{Chunking Topologies: Fixed slicing breaks mid-sentence; recursive splitting honors document separators; semantic chunking splits when embedding similarity drops.}
\label{fig:v2-ch02-chunking}
\end{figure}

\subsection{Recursive Character Splitting (Done Right)}
Recursive splitting tries a prioritized list of natural boundaries—paragraphs (\texttt{"\textbackslash n\textbackslash n"}), lines (\texttt{"\textbackslash n"}), sentences (\texttt{". "}), clauses (\texttt{"; "}), words (\texttt{" "}), and finally raw characters (\texttt{""})—as illustrated in Figure~\ref{fig:v2-ch02-chunking}B.

A classic bug in naive implementations is computing \texttt{new\_separators = separators[i+1:]} and forgetting to \emph{recurse} when an oversized paragraph is encountered, or forgetting to include \texttt{len(separator)} when budgeting overlap. Listing~\ref{lst:v2-ch02-recursive} fixes both bugs cleanly.

\begin{lstlisting}[language=Python, label={lst:v2-ch02-recursive}, caption={Production Recursive Separator Splitter with True Recursion and Exact Overlap}]
import re

class ProductionRecursiveSplitter:
    def __init__(self, chunk_size: int = 1000, chunk_overlap: int = 150):
        self.chunk_size = chunk_size
        self.chunk_overlap = chunk_overlap
        self.separators = ["\n\n", "\n", ". ", "; ", ", ", " ", ""]

    def split_text(self, text: str) -> list[str]:
        def _split(txt: str, separators: list[str]) -> list[str]:
            if len(txt) <= self.chunk_size:
                return [txt] if txt.strip() else []

            separator = separators[-1]
            new_separators: list[str] = []
            for i, sep in enumerate(separators):
                if sep == "":
                    separator = ""
                    break
                if re.search(re.escape(sep), txt):
                    separator = sep
                    new_separators = separators[i + 1:]
                    break

            splits = txt.split(separator) if separator else list(txt)
            docs: list[str] = []
            current_doc: list[str] = []
            current_len = 0

            def _flush() -> None:
                nonlocal current_doc, current_len
                if not current_doc:
                    return
                merged = separator.join(current_doc)
                if merged.strip():
                    docs.append(merged)
                overlap_doc: list[str] = []
                overlap_len = 0
                for prev_s in reversed(current_doc):
                    added = len(prev_s) + (len(separator) if overlap_doc else 0)
                    if overlap_len + added > self.chunk_overlap:
                        break
                    overlap_doc.insert(0, prev_s)
                    overlap_len += added
                current_doc = overlap_doc
                current_len = overlap_len

            for s in splits:
                if len(s) > self.chunk_size:
                    _flush()
                    current_doc, current_len = [], 0
                    if new_separators:
                        docs.extend(_split(s, new_separators))
                    else:
                        step = max(1, self.chunk_size - self.chunk_overlap)
                        for start in range(0, len(s), step):
                            part = s[start:start + self.chunk_size]
                            if part.strip():
                                docs.append(part)
                    continue

                sep_len = len(separator) if current_doc else 0
                if current_len + sep_len + len(s) > self.chunk_size:
                    _flush()
                    sep_len = len(separator) if current_doc else 0
                    if current_len + sep_len + len(s) > self.chunk_size:
                        current_doc, current_len, sep_len = [], 0, 0

                current_doc.append(s)
                current_len += sep_len + len(s)

            if current_doc:
                merged = separator.join(current_doc)
                if merged.strip():
                    docs.append(merged)
            return docs

        return _split(text, self.separators)
\end{lstlisting}

\textbf{Step-by-Step Walkthrough on a Concrete Input:} Configure \texttt{chunk\_size=80, chunk\_overlap=20} and pass:
\begin{quote}\small\ttfamily
"Section 1: Overview. Maintain liquidity.\textbackslash n\textbackslash nSection 2: Covenants. EBITDA must exceed \$50M. The debt ratio shall not exceed 3.5:1.0.\textbackslash n\textbackslash nSection 3: Remedies."
\end{quote}
\begin{enumerate}[leftmargin=18pt]
  \item Line 11 checks \texttt{len(txt) <= 80} (False, 177 chars). Lines 16--23 match the highest-priority separator \texttt{"\textbackslash n\textbackslash n"} and set \texttt{new\_separators = ["\textbackslash n", ". ", "; ", ", ", " ", ""]}.
  \item Line 25 splits into three paragraphs: $p_0$ (42 chars), $p_1$ (92 chars), and $p_2$ (20 chars).
  \item $p_0$ (42 chars $\le 80$) enters \texttt{current\_doc}. Next comes $p_1$ (92 chars $> 80$). Line 49 catches \texttt{len(s) > chunk\_size}, flushes $p_0$ as Chunk 0, clears the buffer, and calls \texttt{\_split(p\_1, new\_separators)} at line 53!
  \item Inside the recursive call on $p_1$, \texttt{"\textbackslash n"} is absent, so line 20 selects \texttt{". "} and splits into three sentences: \texttt{"Section 2: Covenants"} (20), \texttt{"EBITDA must exceed \$50M"} (23), and \texttt{"The debt ratio shall not exceed 3.5:1.0."} (40). The first two join with \texttt{". "} to 45 chars; adding the third ($45 + 2 + 40 = 87 > 80$) triggers \texttt{\_flush()}. In \texttt{\_flush()}, \texttt{overlap\_doc} retains \texttt{"EBITDA must exceed \$50M"} (23 chars $> 20$? Wait: $23 > 20$, so overlap stops at 0 if overlap=20, or retains it if \texttt{chunk\_overlap=25}—and crucially line 40 accounts for \texttt{len(separator)} whenever two retained segments are joined!).
  \item Every emitted chunk is guaranteed to obey \texttt{len(chunk) <= 80} while preserving sentence boundaries.
\end{enumerate}

\subsection{Semantic Chunking: A Worked 4-Sentence Example}
Even recursive splitting cuts blindly between paragraphs that discuss the same topic, or merges two unrelated paragraphs that happen to fit under \texttt{chunk\_size}. \textbf{Semantic chunking} splits when the meaning shifts between consecutive sentences.

Suppose a runbook contains 4 sentences:
\begin{itemize}[leftmargin=18pt]
  \item $s_0$: \emph{``PostgreSQL records every data change in the Write-Ahead Log before modifying heap pages.''}
  \item $s_1$: \emph{``If the server loses power, WAL replay restores committed transactions safely.''}
  \item $s_2$: \emph{``Kubernetes Horizontal Pod Autoscalers monitor CPU utilization every 15 seconds.''}
  \item $s_3$: \emph{``When load exceeds the target threshold, the controller adds new pod replicas.''}
\end{itemize}
Suppose we embed each sentence (or 3-sentence sliding window) and compute cosine similarity between consecutive pairs using Chapter~1's formula, obtaining:
\[
S_{\cos}(s_0, s_1) = 0.88, \qquad S_{\cos}(s_1, s_2) = 0.42, \qquad S_{\cos}(s_2, s_3) = 0.85
\]
Converting similarities to cosine distances $D_i = 1 - S_{\cos}(s_i, s_{i+1})$ gives:
\[
D = [D_0, D_1, D_2] = [0.12,\; 0.58,\; 0.15]
\]
Compute the statistics of $D$:
\[
\mu_D = \frac{0.12 + 0.58 + 0.15}{3} \approx 0.2833, \qquad \sigma_D = \sqrt{\frac{(0.12-0.2833)^2 + (0.58-0.2833)^2 + (0.15-0.2833)^2}{3}} \approx 0.2101
\]
Using an 85th-percentile cutoff (linearly interpolated on sorted $[0.12, 0.15, 0.58]$ at index $0.85 \times 2 = 1.7$, giving $\text{cutoff} = 0.15 + 0.7(0.58 - 0.15) = 0.451$), or equivalently $\tau = \mu_D + 0.8\sigma_D \approx 0.4514$:
\begin{itemize}[leftmargin=18pt]
  \item $D_0 = 0.12 \le 0.451 \implies$ keep $s_0$ and $s_1$ together.
  \item $D_1 = 0.58 > 0.451 \implies$ \textbf{topic jump detected!} Split between $s_1$ and $s_2$.
  \item $D_2 = 0.15 \le 0.451 \implies$ keep $s_2$ and $s_3$ together.
\end{itemize}
Result: Chunk 1 = $\{s_0, s_1\}$ (PostgreSQL WAL) and Chunk 2 = $\{s_2, s_3\}$ (Kubernetes HPA). Listing~\ref{lst:v2-ch02-semantic} implements this exact pass.

\begin{lstlisting}[language=Python, label={lst:v2-ch02-semantic}, caption={Semantic Chunking via Adjacent Cosine Distance}]
import numpy as np

def semantic_chunking(
    sentences: list[str],
    embeddings: np.ndarray,
    threshold_percentile: float = 85.0
) -> list[str]:
    """Splits sentences into chunks where embedding distance spikes."""
    if len(sentences) <= 1:
        return sentences

    # embeddings shape: (N, D), pre-normalized to unit length
    distances = np.array([
        1.0 - float(np.dot(embeddings[i], embeddings[i + 1]))
        for i in range(len(embeddings) - 1)
    ])
    cutoff = float(np.percentile(distances, threshold_percentile))

    chunks: list[str] = []
    current_chunk = [sentences[0]]
    for i, dist in enumerate(distances):
        if dist > cutoff:
            chunks.append(" ".join(current_chunk))
            current_chunk = [sentences[i + 1]]
        else:
            current_chunk.append(sentences[i + 1])

    if current_chunk:
        chunks.append(" ".join(current_chunk))
    return chunks
\end{lstlisting}

\section{The ``Lost in the Middle'' Phenomenon}

Even if chunking is spotless, pasting 30 retrieved chunks into a 64k-token prompt backfires. Liu et al. (2023) showed that LLM accuracy forms a U-shaped curve (Figure~\ref{fig:v2-ch02-lost-middle}): models reliably use facts at the very start (\textbf{primacy}) and very end (\textbf{recency}) of the prompt, while accuracy drops by 30--50\% when the critical answer sits in the middle.

\begin{figure}[htbp]
\centering
\begin{tikzpicture}
\draw[->, thick, color=headingblue] (0,0) -- (10,0) node[right, font=\sffamily\footnotesize] {Relative Position in Prompt Context};
\draw[->, thick, color=headingblue] (0,0) -- (0,4.3) node[above, font=\sffamily\footnotesize] {QA Accuracy (\%)};

\draw[very thick, color=warningRed, domain=0.5:9.5, samples=100]
  plot (\x, {0.08*(\x - 5.0)*(\x - 5.0) + 1.2});

\node[font=\sffamily\scriptsize\bfseries, color=linkblue] at (1.8, 3.7) {Primacy Effect (Start)};
\node[font=\sffamily\scriptsize\bfseries, color=linkblue] at (8.2, 3.7) {Recency Effect (End)};
\node[font=\sffamily\scriptsize\bfseries, color=warningRed] at (5.0, 0.75) {Middle Drop ($-30\%$ to $-50\%$)};

\draw[dashed, color=gray!50] (0.5, 3.2) -- (0.5, 0) node[below, font=\tiny\sffamily] {0\%};
\draw[dashed, color=gray!50] (5.0, 1.2) -- (5.0, 0) node[below, font=\tiny\sffamily] {50\%};
\draw[dashed, color=gray!50] (9.5, 3.2) -- (9.5, 0) node[below, font=\tiny\sffamily] {100\%};

\end{tikzpicture}
\caption{The ``Lost in the Middle'' U-Shaped Attention Curve: Information buried near the center of a long prompt receives weaker attention than chunks placed near the start or end.}
\label{fig:v2-ch02-lost-middle}
\end{figure}

\begin{underthehoodbox}[Why Does This Happen? Attention Decay, RoPE, and Attention Sinks]
When an LLM generates the next word, its \textbf{self-attention} mechanism assigns a numerical weight (summing to 1.0) across all earlier tokens in the prompt to decide which ones matter right now. Two mechanics starve middle tokens:
\begin{enumerate}
  \item \textbf{Rotary Position Embeddings (RoPE) Distance Decay}: To tell whether one word is 5 tokens away or 20,000 tokens away, modern LLMs encode position by rotating each token's internal vector in 2D planes by an angle proportional to its position index. As the gap $|i - j|$ between a middle chunk and the end of the prompt stretches across tens of thousands of tokens, relative alignment decays and long-range attention blurs.
  \item \textbf{Attention Sinks and Recency}: Tokens at the very end sit immediately next to the freshly generated output, dominating recency weights. Meanwhile, the very first tokens (token 0, the system prompt) act as an \textbf{attention sink}: because attention weights must sum to 1.0 even when no specific word is strongly relevant, models dump spare attention weight onto the opening tokens, leaving middle tokens starved.
\end{enumerate}
\end{underthehoodbox}

\begin{productionwarningbox}[The Context-Stuffing Anti-Pattern]
Upgrading from an 8k-token model to a 128k-token model does not fix poor retrieval. Dumping 40 raw chunks into a prompt buries the best evidence in the middle and inflates latency and cost. Keep retrieved context tight (top 5--8 chunks after reranking, Chapter~3), deduplicate overlapping chunks, and place your highest-scoring chunks at the top and bottom of the context block.
\end{productionwarningbox}

\section{Small-to-Big Retrieval and Contextual Enrichment}

\subsection{Parent-Document Retrieval: Decoupling Search from Generation}
Chunk size creates a tug-of-war:
\begin{itemize}[leftmargin=18pt]
  \item \textbf{Small chunks (100--150 tokens)} produce razor-sharp vectors because each vector describes one focused fact.
  \item \textbf{Large chunks (1,000--1,500 tokens)} give the LLM full surrounding clauses, definitions, and tables.
\end{itemize}
\textbf{Parent-Document Retrieval} solves both by indexing small \emph{child} chunks in the vector database while storing the large \emph{parent} section in a key-value store (Figure~\ref{fig:v2-ch02-parent}).

\begin{figure}[htbp]
\centering
\begin{tikzpicture}[
  parent/.style={draw=headingblue, fill=codebg, rounded corners=2pt, minimum width=9.0cm, minimum height=1.4cm, font=\sffamily\scriptsize, align=center},
  child/.style={draw=exampleBlue, fill=exampleBg, rounded corners=1pt, minimum width=2.4cm, minimum height=0.75cm, font=\sffamily\tiny\bfseries, align=center},
  arrow/.style={-{Stealth[scale=0.8]}, thick, color=headingblue}
]

\node[parent] (p1) at (5, 2.4) {Parent Section (1,500 tokens, \texttt{parent\_id="p-42"})\\Stored in Document Store (PostgreSQL / Redis)};
\node[child] (c1) at (1.5, 0.7) {Child 1 (150 tok)\\Vector $\vec{v}_1$};
\node[child] (c2) at (5.0, 0.7) {Child 2 (150 tok)\\Vector $\vec{v}_2$};
\node[child] (c3) at (8.5, 0.7) {Child 3 (150 tok)\\Vector $\vec{v}_3$};

\draw[arrow] (p1.south) -| (c1.north);
\draw[arrow] (p1.south) -- (c2.north);
\draw[arrow] (p1.south) -| (c3.north);

\node[draw=intuitionTeal, fill=intuitionBg, rounded corners=2pt, font=\sffamily\scriptsize\bfseries, minimum width=9.0cm, minimum height=0.75cm] at (5, -0.65) {
  Vector search matches Child 2 $\implies$ Look up \texttt{p-42} $\implies$ Return full Parent Section to LLM
};

\end{tikzpicture}
\caption{Parent-Document (Small-to-Big) Retrieval: Granular child vectors maximize search accuracy; deduplicated parent sections preserve full context for generation.}
\label{fig:v2-ch02-parent}
\end{figure}

\begin{lstlisting}[language=Python, label={lst:v2-ch02-parent}, caption={Parent-Document Retrieval with Ordered Parent Deduplication}]
import uuid

class ParentDocumentRetriever:
    def __init__(self, vector_store, doc_store: dict[str, str]):
        self.vector_store = vector_store
        self.doc_store = doc_store
        self.child_splitter = ProductionRecursiveSplitter(chunk_size=200, chunk_overlap=20)
        self.parent_splitter = ProductionRecursiveSplitter(chunk_size=1500, chunk_overlap=150)

    def add_documents(self, raw_documents: list[str]) -> None:
        for raw_doc in raw_documents:
            parents = self.parent_splitter.split_text(raw_doc)
            for parent_text in parents:
                parent_id = str(uuid.uuid4())
                self.doc_store[parent_id] = parent_text

                children = self.child_splitter.split_text(parent_text)
                for child_text in children:
                    metadata = {"parent_id": parent_id, "snippet": child_text[:50]}
                    self.vector_store.add_text(child_text, metadata=metadata)

    def retrieve(self, query: str, top_k_children: int = 10) -> list[str]:
        child_matches = self.vector_store.search(query, k=top_k_children)

        seen_parent_ids: set[str] = set()
        parent_documents: list[str] = []
        for match in child_matches:
            pid = match.metadata["parent_id"]
            if pid not in seen_parent_ids:
                seen_parent_ids.add(pid)
                parent_documents.append(self.doc_store[pid])
        return parent_documents
\end{lstlisting}

\textbf{Step-by-Step Trace of \texttt{ParentDocumentRetriever}:}
Suppose we ingest our bond indenture from the motivating problem:
\begin{enumerate}[leftmargin=18pt]
  \item \texttt{add\_documents} splits the indenture into 1,500-character parents (lines 12--15). Section 4.02 (Negative Covenants) gets \texttt{parent\_id = "p-402"} and its full text is saved in \texttt{self.doc\_store["p-402"]}.
  \item Lines 17--20 slice \texttt{"p-402"} into 200-character children: Child A mentions \emph{``Section 4.02 Negative Covenants: Consolidated EBITDA''}, Child B contains \emph{``...the ratio shall not exceed 3.5:1.0''}. Both carry \texttt{metadata=\{"parent\_id": "p-402"\}}.
  \item When the analyst queries \texttt{"debt ratio under negative covenants"}, line 23 returns Child B (rank 1, \texttt{parent\_id="p-402"}) and Child A (rank 2, \texttt{parent\_id="p-402"}), followed by Child C from Section 7.1 (\texttt{parent\_id="p-701"}).
  \item Lines 27--31 iterate in rank order: \texttt{"p-402"} is added once to \texttt{seen\_parent\_ids} and appended; when Child A arrives with the same \texttt{"p-402"}, line 29 skips the duplicate so the LLM receives Section 4.02 once—complete with its covenant header—followed by Section 7.1.
\end{enumerate}

\subsection{Contextual Retrieval: Eliminating Orphan Chunks}
Consider a standalone chunk from an SEC filing:
\begin{center}
\emph{``Operating margin declined by 2.4\% due to regulatory headwinds in the mobility segment.''}
\end{center}
Even with small chunks, a query asking \emph{``Why did Uber's Q3 2024 operating margin drop?''} fails if neither \emph{``Uber''} nor \emph{``Q3 2024''} appears inside the chunk! \textbf{Contextual Retrieval} (Anthropic, 2024) fixes this at ingestion time by asking a fast, low-cost LLM to prepend a 2-sentence situating prefix before embedding (Listing~\ref{lst:v2-ch02-contextual}).

\begin{lstlisting}[language=Python, label={lst:v2-ch02-contextual}, caption={Contextual Enrichment Prior to Vector Embedding}]
CONTEXT_ENRICHMENT_PROMPT = """<document>
{full_document}
</document>

Here is the chunk we want to situate within the whole document:
<chunk>
{chunk_text}
</chunk>

Give a short, concise context (2-3 sentences) to situate this chunk within the overall document for search retrieval. Answer ONLY with the context."""

async def generate_contextual_chunk(
    llm_client, full_document: str, chunk_text: str
) -> str:
    prompt = CONTEXT_ENRICHMENT_PROMPT.format(
        full_document=full_document[:20000],
        chunk_text=chunk_text
    )
    context_prefix = await llm_client.generate(prompt, max_tokens=100)
    return f"{context_prefix.strip()}\n\n{chunk_text}"
\end{lstlisting}

Prepending \emph{``This chunk is from Uber Technologies' Q3 2024 Form 10-Q discussing Mobility segment results.''} embeds \texttt{Uber}, \texttt{Q3 2024}, and \texttt{Form 10-Q} directly into the chunk vector—cutting retrieval failures by nearly half in production benchmarks, especially when paired with prompt caching (Chapter~7) so \texttt{full\_document} is cached across all chunks of the same file.

\begin{realworldbox}[Diagnosing Orphaned Pronoun \& Table-Header Failures]
An internal compliance assistant answered \emph{``I found no record of the Q3 margin change''} on 38\% of financial questions. Inspecting the retrieval trace exposed the culprit:
\begin{lstlisting}[basicstyle=\ttfamily\scriptsize]
TRACE query="What was Acme Corp's Q3 2024 operating margin?"
  top_1_hit (cos=0.61): "2024 | $150M | $95M | $55M"         # Missing header \& company!
  top_2_hit (cos=0.58): "The segment fell 2.4% vs prior yr." # Missing entity \& quarter!
\end{lstlisting}
The true table row had been split away from its header, and the narrative paragraph never repeated \emph{``Acme Corp''}. Flattening table rows into prose tuples (Listing~\ref{lst:v2-ch02-tables}) and prepending contextual prefixes (Listing~\ref{lst:v2-ch02-contextual}) lifted top-5 retrieval recall from 62\% to 94\%.
\end{realworldbox}

\begin{tryityourselfbox}[Tracing Recursive Splitting \& Overlap Calculation]
Trace \texttt{ProductionRecursiveSplitter(chunk\_size=50, chunk\_overlap=15)} on the input:
\begin{quote}\small\ttfamily
"Alpha one. Beta two. Gamma three.\textbackslash n\textbackslash nDelta four."
\end{quote}
List the exact output chunks, their character lengths, and explain which separator level handled each split. (Worked solution is in Section~2.7.)
\end{tryityourselfbox}

\section{Chapter Summary \protect\& Key Takeaways}
\begin{itemize}[leftmargin=18pt]
  \item \textbf{Structure-First Ingestion}: Use layout-aware parsers for multi-column PDFs and flatten tables into self-contained propositions so rows keep their headers.
  \item \textbf{Recursive \& Semantic Chunking}: True recursive splitters descend into finer separators whenever a segment exceeds \texttt{chunk\_size} and account for \texttt{len(separator)} during overlap; semantic chunkers detect topic shifts via spikes in adjacent cosine distance.
  \item \textbf{Lost in the Middle}: Due to RoPE distance decay, recency, and attention sinks, LLMs under-attend to the center of long prompts. Rerank and place top chunks at the edges.
  \item \textbf{Parent-Document \& Contextual Retrieval}: Search over small, focused child chunks, deduplicate \texttt{parent\_id}s, and return full parent sections; prepend document/section context before embedding to eliminate orphan pronouns.
\end{itemize}

\section{Exercises}

\subsection*{Mechanics}
\textbf{Exercise \thechapter.1}: A document has 10 sentences with adjacent cosine distances:
\[
D = [0.12,\; 0.15,\; 0.08,\; 0.45,\; 0.10,\; 0.11,\; 0.52,\; 0.09,\; 0.14]
\]
Given $\tau = \mu_D + 1.0 \cdot \sigma_D$, compute $\mu_D$, $\sigma_D$, and identify which sentence boundaries trigger splits.

\textbf{Exercise \thechapter.2}: Explain the difference between character overlap and token overlap. If a splitter sets \texttt{chunk\_overlap=100} characters on Python source code, approximately how many tokens of overlap does that produce?

\subsection*{Debugging}
\textbf{Exercise \thechapter.3}: After deploying \texttt{ParentDocumentRetriever}, an engineer notices that 5 out of 5 returned parent documents in the prompt are identical copies of the same 1,500-token section, crowding out all other evidence. Explain why this happens and write the Python deduplication fix.

\subsection*{Architecture}
\textbf{Exercise \thechapter.4}: Design an incremental ingestion pipeline for an enterprise wiki (Notion/Confluence) that updates parent and child records when a single page is edited, without leaving orphaned child vectors in the index.

\section{Detailed Solutions}

\subsection*{Solution \thechapter.1}
Count $N = 9$, sum $= 1.76$, so $\mu_D = 1.76 / 9 \approx 0.1956$. Computing squared deviations $(D_i - 0.1956)^2$:
\[
0.0057 + 0.0021 + 0.0134 + 0.0647 + 0.0091 + 0.0073 + 0.1052 + 0.0112 + 0.0031 = 0.2218
\]
Variance $\sigma_D^2 = 0.2218 / 9 \approx 0.02464 \implies \sigma_D \approx 0.1570$. Threshold $\tau = 0.1956 + 0.1570 = 0.3526$. Two distances exceed $\tau$: $D_3 = 0.45$ (split between $s_3$ and $s_4$) and $D_6 = 0.52$ (split between $s_6$ and $s_7$), yielding $\{s_0..s_3\}$, $\{s_4..s_6\}$, and $\{s_7..s_9\}$.

\subsection*{Solution \thechapter.2}
Character overlap counts raw characters (`len(s)`), whereas token overlap counts BPE subword pieces. In English prose, $1\text{ token} \approx 4\text{ chars}$; in code (punctuation, indentation, short identifiers), $1\text{ token} \approx 2.5\text{ to }3.0\text{ chars}$. Thus 100 characters of Python yields $\approx 100 / 2.8 \approx 35\text{--}40$ tokens.

\subsection*{Solution \thechapter.3}
Multiple child chunks from the same parent document ranked in the top 5 vector hits, each carrying the same \texttt{parent\_id}. Fetching \texttt{doc\_store[pid]} without tracking seen IDs repeats the parent 5 times. Fix:
\begin{lstlisting}[language=Python]
def deduplicate_parents(child_hits: list, doc_store: dict[str, str]) -> list[str]:
    seen: set[str] = set()
    parents: list[str] = []
    for hit in child_hits:
        pid = hit.metadata["parent_id"]
        if pid not in seen:
            seen.add(pid)
            parents.append(doc_store[pid])
    return parents
\end{lstlisting}

\subsection*{Solution \thechapter.4}
Key each parent and child deterministically by \texttt{sha256(doc\_id + ":" + str(version) + ":" + str(idx))}. On page update: (1) write the new parent and child vectors tagged with \texttt{doc\_id} and \texttt{version = v\_new}; (2) atomically update the pointer/metadata; (3) delete stale chunks with \texttt{DELETE FROM chunks WHERE doc\_id = :id AND version < :v\_new}; (4) route transient failures to a Dead-Letter Queue for idempotent retry.

\subsection*{Solution (Try It Yourself)}
First, top-level \texttt{\_split} splits on \texttt{"\textbackslash n\textbackslash n"} into $p_0 = \text{"Alpha one. Beta two. Gamma three."}$ (33 chars $\le 50$) and $p_1 = \text{"Delta four."}$ (11 chars $\le 50$). Adding $p_1$ to $p_0$ would reach $33 + 2 + 11 = 46 \le 50$, so wait: \texttt{"\textbackslash n\textbackslash n".join([p0, p1])} has length 46, fitting in a single 46-char chunk! If instead $p_0$ were 45 chars and $p_1$ were 11 chars ($45 + 2 + 11 = 58 > 50$), \texttt{\_flush()} emits $p_0$, checks overlap ($45 > 15$, so overlap is empty), and emits $p_1$ separately. (Testing edge cases like this with \texttt{len(separator)} included is why exact budget accounting matters!)

\begin{labbridgebox}[Hands-On Exercise \& Project Bridge]
Review how chunked documents feed dense cosine retrieval in \texttt{exercises/volume2/ex01\_vector\_math\_cosine.py}, and explore how retrieved text segments enter the hybrid search pipeline in \texttt{projects/p06-rag-hybrid-rerank/}.
\end{labbridgebox}

\section{Further Reading \protect\& Citations}
\begin{itemize}[leftmargin=18pt]
  \item Liu, N. F., et al. (2023). \emph{Lost in the Middle: How Language Models Use Long Contexts}. TACL.
  \item Anthropic. (2024). \emph{Introducing Contextual Retrieval}. Anthropic Research.
  \item Lewis, P., et al. (2020). \emph{Retrieval-Augmented Generation for Knowledge-Intensive NLP Tasks}. NeurIPS.
\end{itemize}
